
\subsubsection*{5.2.4}
\phantomsection
\label{sga3:VIIB:5.2.4}

For every pseudocompact vector space \(V\) over \(k\), denote by
\[
  {}_{\pi}V
  =
  V\widehat\otimes_{\pi}k
\]
the space obtained from \(V\) by extension of scalars along
\(x\mapsto x^p\).
We then have a commutative diagram with exact rows:

\begin{figure}[H]
  \centering
  \begin{tikzcd}[row sep=2.8em,column sep=2.6em]
0 \arrow[r]
  & {}_\pi\overline{I^2} \arrow[r,"\alpha"] \arrow[d,"u"']
  & {}_\pi I \arrow[r,"\beta"] \arrow[d,"v"']
  & {}_\pi\omega \arrow[r] \arrow[d,"w"']
  & 0 \\
0 \arrow[r]
  & \overline{J^2} \arrow[r,"\gamma"]
  & J \arrow[r,"\delta"]
  & \overline\omega \arrow[r]
  & 0\, ,
\end{tikzcd}

\end{figure}

\noindent
Here
\[
  \overline\omega=J/J^2,
\]
and the maps \(u,v,w\) are induced by the linear map
\[
  {}_{\pi}B\longrightarrow B^p,
  \qquad
  x\widehat\otimes a\longmapsto x^pa.
\]
Since \(u\) and \(v\) are surjective, so is \(w\), and its kernel
is the image \({}_{\pi}\omega_1\) of
\[
  {}_{\pi}I_1=\ker(v).
\]

Put
\[
  J_n=\{\,x\in J\mid x^{p^n}=0\,\},
  \qquad
  \overline\omega_n=\delta(J_n).
\]
Then, for every \(n\geq0\),
\[
  J_n=v({}_{\pi}I_{n+1}),
  \qquad
  \overline\omega_n=w({}_{\pi}\omega_{n+1}).
\]

\phantomsection
\label{sga3:VIIB:5.2.4:delta-source-note}
\noindent\textit{Translator's note.}\enspace The French text prints
\(\overline\omega_n=\delta(I_n)\), although the displayed exact
row defines \(\delta\) on \(J\), and the immediately preceding
definition introduces \(J_n\subset J\). The type-correct
\(\delta(J_n)\) is used above; all subsequent identities agree
with this reading.

The section
\[
  {}_{\pi}\sigma:{}_{\pi}\omega\longrightarrow{}_{\pi}I
\]
induced by the section \(\sigma\) of \SGARefLink{sga3:VIIB:5.2.1:section-proposition}{5.2.1} therefore descends to
a section
\[
  \tau:\overline\omega\longrightarrow J
\]
compatible with the filtrations of \(J\) and
\(\overline\omega\). Since \(B^p\) has height at most \(r-1\),
the induction hypothesis shows that this section induces an
isomorphism
\[
  \psi:
  B''
  =
  k[[\overline\omega]]
  \big/
  \bigl((x^{p^n})\bigr)_{x\in\overline\omega_n}
  \xrightarrow{\ \sim\ }
  B^p.
\]
But \(B''\) identifies with \(B'^p\), and \(\psi\) identifies
with the morphism \(B'^p\to B^p\) induced by \(\phi\). Thus the
theorem is proved when \(B\) has finite height.
